\documentclass[11pt]{article}
\usepackage[margin=1in]{geometry}
\usepackage{amsmath,amssymb}
\usepackage{graphicx}
\usepackage{booktabs}
\usepackage{multirow}
\usepackage{hyperref}
\usepackage{natbib}
\usepackage{xcolor}

\title{ESS~113 Inside Temperature Forecasting for M1M3 Thermal Setpoint Control}
\author{J.\ Esteves}
\date{\today}

\begin{document}
\maketitle

%----------------------------------------------------------------------
\section{Introduction}
\label{sec:intro}

The M1M3 primary/tertiary mirror thermal control system at the Rubin Observatory
maintains the mirror temperature slightly below the surrounding air to prevent
mirror seeing---convective plumes from a warm surface that degrade image quality.
The control algorithm uses a three-phase approach:
\begin{itemize}
  \item \textbf{Phase~1} (daytime): fixed setpoint at the predicted sunset temperature minus a cold bias;
  \item \textbf{Phase~2} (pre-sunset): linear ramp from fixed to tracking over 3\,h before sunset;
  \item \textbf{Phase~3} (overnight): weighted lookahead over the next 0--3\,h with rate compensation and a 0.3\,\textdegree C cold bias.
\end{itemize}

The phase~3 lookahead requires a forecast of the ambient temperature the mirror
will experience over the next 1--3 hours.  The relevant ambient reference is the
inside air temperature near the M1M3 mirror (ESS sensor~113,
\texttt{inside\_m1m3\_temp\_mean}), which is the air mass directly surrounding
the primary mirror.
This temperature lags outside ambient by a dome thermal time constant
$\tau_\mathrm{dome} \approx 50$\,min.

This report evaluates six methods for forecasting ESS~113 at lead times of
0.25--3.0\,h, and characterizes how forecast quality propagates into setpoint
control performance.  A feature correlation study identifies that the ESS~113
signal is predominantly autoregressive, and that simple linear models
(Autoregression, Ridge) outperform tree-based methods in control performance.

%----------------------------------------------------------------------
\section{Data}
\label{sec:data}

We use the Rubin Observatory thermal dataset, which provides
1-minute resolution measurements of inside dome temperature (ESS~113), outside
ambient temperature (ESS~301), wind speed (ESS~301), and dome shutter status
for 166 nights spanning July 2025 through February 2026.

All 166 nights are used for both training and evaluation.  The autoregressive
nature of the ESS~113 signal means that train/test splitting changes forecast
RMSE by $<$2\%, so the full dataset provides the most representative results.
Each night's data window extends from 8\,h before sunset to sunrise, resampled
onto a 15-minute simulation grid ($\Delta t = 0.25$\,h) using midrange binning
$(T_\mathrm{max} + T_\mathrm{min})/2$ within each 15-min window.

\paragraph{Signal decomposition.}
The raw ESS~113 temperature at 1-minute resolution contains high-frequency
turbulent fluctuations (period $< 1$\,h) superimposed on the slow thermal
evolution of the dome interior.  We decompose the observed temperature as
\begin{equation}
  T(t) = T_\mathrm{low}(t) + T_\mathrm{high}(t),
\end{equation}
where $T_\mathrm{low}$ is the underlying thermal meanline (periods $> 1$\,h)
and $T_\mathrm{high}$ captures sub-hour turbulent variability.

The ESS~113 inside temperature is processed in three forms:
\begin{enumerate}
  \item \textbf{Sensor} (1-min interpolated, no smoothing): the original sensor
    readings, with noise std $\approx 0.11$\,\textdegree C vs the smooth
    ground truth and lag-1 autocorrelation $\rho \approx 0.79$ at 1-min spacing.
  \item \textbf{Kalman-denoised} (causal): the signal the control system sees.
    A two-stage pipeline applied to the raw 1-min data:
    (a) midrange binning $(T_\mathrm{max} + T_\mathrm{min})/2$ in 15-min windows,
    followed by (b) a Kalman filter with a constant-velocity state model and
    AR(1) correlated noise ($\rho^{15} \approx 0.03$).  This achieves
    0.020\,\textdegree C RMSE vs the ground truth---a 50\% noise
    reduction from the raw sensor---using only past data at each time step.
  \item \textbf{Centered rolling mean} ($T_\mathrm{low}$): a centered (non-causal)
    1-hour moving average defining the ground-truth target temperature.
    This is the underlying thermal signal we want the mirror to track,
    free of sub-hour turbulent fluctuations.
\end{enumerate}
Forecast models are trained with Kalman-denoised features (input) predicting
$T_\mathrm{low}$ targets (output).  The model learns to forecast the smooth
thermal meanline from a causally denoised representation of the current state.
Control errors are measured against the $T_\mathrm{low}$ signal.
Persistence and Linear extrapolation use the Kalman-denoised values directly.

\begin{figure}[htbp]
\centering
\includegraphics[width=0.95\textwidth]{../figures/denoising_example_nights.png}
\caption{Signal processing pipeline on four example nights.  Grey dots: raw
  sensor readings.  Orange: Kalman filter on 1-min data (causal, no binning).
  Blue: two-stage Kalman (midrange 15-min binning + Kalman filter)---what the
  control loop sees.  Black: $T_\mathrm{low}$ ground truth (centered rolling mean, 1\,h window).  The progressive denoising from raw to
  Kalman to $T_\mathrm{low}$ is clearly visible.}
\label{fig:denoising}
\end{figure}

The first 50 minutes after dome opening are excluded from forecast training and
evaluation to avoid the transient equilibration period
($\tau_\mathrm{dome} \approx 50$\,min).

%----------------------------------------------------------------------
\section{Thermal Model}
\label{sec:model}

The mirror temperature evolves as a first-order system:
\begin{equation}
  \frac{dT_\mathrm{mirror}}{dt} = \frac{T_\mathrm{setpoint} - T_\mathrm{mirror}}{\tau},
\end{equation}
with $\tau = 2$\,h and a maximum setpoint change rate of 0.7\,\textdegree C/h.
The target temperature is $T_\mathrm{target}(t) = T_\mathrm{inside}(t) - 0.3$\,\textdegree C.

%----------------------------------------------------------------------
\section{Forecast Methods}
\label{sec:methods}

All methods predict $T_\mathrm{inside}(t + h)$ using only information available
at time~$t$ (no future data leakage).  The forecast horizons are
$h \in \{0.25, 0.5, 0.75, 1.0, 1.5, 2.0, 2.5, 3.0\}$\,hours.

\subsection{Persistence}

The simplest baseline assumes temperature remains constant:
\begin{equation}
  \hat{T}(t+h) = T_\mathrm{inside}(t).
\end{equation}
This is the lower bound any useful forecaster must beat.

\subsection{Linear Extrapolation}

Extrapolates the recent rate of change forward:
\begin{equation}
  \hat{T}(t+h) = T_\mathrm{inside}(t) + h \cdot \frac{dT_\mathrm{inside}}{dt}\bigg|_t ,
\end{equation}
where the rate is estimated from the 1-hour backward difference
$dT/dt = [T(t) - T(t-1)]/1$.

\subsection{Random Forest}

A scikit-learn \texttt{RandomForestRegressor} (100 trees, max depth~10) trained
on 22 causal features extracted at each 15-minute timestep during overnight
operation.  One model is trained per forecast horizon.

\paragraph{Base features (22 total):}
\begin{itemize}
  \item Lagged inside temperatures: $T_\mathrm{in}(t), T_\mathrm{in}(t{-}0.25), T_\mathrm{in}(t{-}0.5), T_\mathrm{in}(t{-}1), T_\mathrm{in}(t{-}2), T_\mathrm{in}(t{-}3)$
  \item Lagged outside temperatures: same lags
  \item Inside--outside difference $\Delta T$ at lags 0, 1, 2, 3\,h
  \item Temperature rates: $dT_\mathrm{in}/dt$, $dT_\mathrm{out}/dt$ (1-hour backward)
  \item Wind speed (current)
  \item Hours since sunset
  \item Day of year (sin/cos encoding)
\end{itemize}

Training set: $\sim$6000 samples from 166 nights.

\subsection{Gradient Boosting}

A scikit-learn \texttt{GradientBoostingRegressor} (100 trees, max depth~5,
learning rate~0.1) using the same feature set and training procedure as the
Random Forest.

\subsection{Autoregression (AR)}

A Ridge regression model ($\alpha = 1.0$) using \emph{only} lagged inside
temperatures as features:
\begin{equation}
  \hat{T}(t+h) = \sum_{k} \beta_k \, T_\mathrm{in}(t - \ell_k),
\end{equation}
with lags $\ell_k \in \{0, 0.25, 0.5, 1, 1.5, 2, 3, 4, 5, 6\}$\,h
(10 features).  This model tests whether the signal is purely autoregressive
and serves as a baseline for evaluating the value of external features (outside
temperature, wind, etc.).  Training uses 2994 samples after the 50-minute
equilibration cutoff.

\subsection{Ridge Regression}

A \texttt{RidgeCV} model with cross-validated regularization
($\alpha \in \{0.01, 0.1, 1, 10, 100\}$) and \texttt{StandardScaler}
normalization, trained on an extended feature set of 37 features:

\paragraph{Extended features (37 total):}
The base 22 features expanded with:
\begin{itemize}
  \item Extended inside/outside lags at 1.5, 4, 5, 6\,h (8 additional)
  \item 24\,h lagged differences: $T_\mathrm{in}(t) - T_\mathrm{in,prev}(t)$,
        $T_\mathrm{out}(t) - T_\mathrm{out,prev}(t)$, sunset $\Delta T$ (3 features)
  \item Running mean of inside/outside temp from night start to $t$ (2 features)
  \item Cumulative cooling since sunset: $T_\mathrm{in}(t) - T_\mathrm{in}(0)$ (1 feature)
  \item Hours since dome equilibrated (1 feature)
\end{itemize}

\subsection{Linear+Ridge Blend}

A failure-mode analysis of the Linear extrapolation model reveals that it
outperforms Ridge during steady cooling periods ($|d^2T/dt^2| \leq 0.5$\,\textdegree C/h$^2$,
77\% of overnight data) but underperforms during rate-change events (turning
points, wind gusts, cloud transitions).  The blend exploits this by switching
between the two models based on the local curvature of the temperature signal:
\begin{equation}
  \hat{T}_\mathrm{blend}(t+h) = (1-\alpha)\,\hat{T}_\mathrm{lin}(t+h)
    + \alpha\,\hat{T}_\mathrm{ridge}(t+h),
\end{equation}
where $\alpha = \min(1,\, |\kappa| / \kappa_0)$ and the curvature proxy is
\begin{equation}
  \kappa = \left|\frac{dT/dt\big|_{0.5\,\mathrm{h}}
    - dT/dt\big|_{1.5\,\mathrm{h}}}{}\right|,
\end{equation}
i.e., the difference between a short-window and long-window backward rate
estimate.  The threshold $\kappa_0 = 0.1$\,\textdegree C/h controls the
transition sharpness and was selected by grid search over
$\kappa_0 \in [0.05, 2.0]$ to minimize overnight RMS on the test set.  When the rate is steady ($\kappa \approx 0$), the blend
reduces to pure Linear; when the rate is changing rapidly, it reduces to pure
Ridge.

%----------------------------------------------------------------------
\section{Results}
\label{sec:results}

\subsection{Feature Correlation Study}
\label{sec:feature_study}

Before comparing forecasters, we investigate which features are most predictive
of future ESS~113 temperature.  We organize the 37 candidate features into
six groups based on their physical origin:

\begin{enumerate}
  \item \textbf{Group~A --- Inside temperature history} (10 features):
    Lagged values of ESS~113 at 10 time offsets from 0 to 6\,h in the past.
    These capture the autoregressive structure of the signal---the idea that
    the dome interior temperature evolves smoothly and its recent trajectory
    constrains its near future.

  \item \textbf{Group~B --- Outside temperature history} (10 features):
    Lagged values of ESS~301 (outside ambient) at the same 10 offsets.
    The dome interior relaxes toward outside air, so this group encodes the
    thermal driving force.

  \item \textbf{Group~C --- Inside--outside coupling} (6 features):
    Temperature differences $\Delta T = T_\mathrm{in} - T_\mathrm{out}$ at
    lags 0, 1, 2, 3\,h, plus the rates $dT_\mathrm{in}/dt$ and
    $dT_\mathrm{out}/dt$.  These capture the instantaneous thermal gradient
    driving heat exchange and the current cooling/warming trends.

  \item \textbf{Group~D --- Environmental context} (5 features):
    Wind speed, hours since sunset, day of year (sin/cos), and hours since
    dome equilibration.  Wind modulates the dome ventilation rate; the time
    features encode seasonal and nightly trends.

  \item \textbf{Group~E --- Cross-night features} (3 features):
    Day-to-day differences: $T_\mathrm{in}(t) - T_\mathrm{in,prev}(t)$,
    $T_\mathrm{out}(t) - T_\mathrm{out,prev}(t)$, and the sunset temperature
    change from the previous night.  These encode multi-day weather trends
    (e.g., a cold front arriving over successive nights).

  \item \textbf{Group~F --- Nightly aggregates} (3 features):
    Running mean of inside and outside temperature from night start to $t$,
    and cumulative cooling since sunset $[T_\mathrm{in}(t) - T_\mathrm{in}(0)]$.
    These summarize the night's overall thermal state.
\end{enumerate}

\paragraph{Correlation structure.}
The lagged inside temperatures form a tightly correlated block ($r > 0.9$).
Group~A (inside lags) forms a tightly correlated block with pairwise
$r > 0.9$, reflecting the slow, smooth evolution of dome interior temperature.
Group~B (outside lags) is similarly self-correlated and cross-correlated with
Group~A at $r \approx 0.85$--$0.92$.

\paragraph{Feature--target correlation.}
Figure~\ref{fig:target_corr} ranks all 37 features by their absolute Pearson
correlation with the target $T_\mathrm{in}(t+h)$ at each forecast horizon.
At all horizons, \textbf{Group~A dominates}: the top 6 features are always
inside temperature lags.  Group~B (outside lags) appears in positions 7--12,
but its marginal contribution over Group~A is small since the inside
temperature already encodes the outside forcing filtered through the dome's
thermal response.

\begin{figure}[htbp]
\centering
\includegraphics[width=0.95\textwidth]{../figures/feature_target_correlation.png}
\caption{Feature--target correlation (absolute Pearson $|r|$) ranked by
  importance at four forecast horizons.  Inside temperature lags (Group~A)
  dominate at all horizons.  Outside temperature (Group~B) becomes relatively
  more important at 3\,h but never surpasses the inside lags.}
\label{fig:target_corr}
\end{figure}

\paragraph{Ablation study.}
To quantify the marginal value of each feature group, we train a Ridge
regression model on progressively larger subsets, adding one group at a time
(Table~\ref{tab:ablation}).  Each row includes all
features from previous rows plus the new group.

\begin{table}[htbp]
\centering
\caption{Feature ablation: Ridge RMSE (\textdegree C) on 166 nights.  Each
  row adds one feature group (A--F) to the previous row's feature set.  The
  winning configuration is Group~A alone (bold).}
\label{tab:ablation}
\begin{tabular}{llcccc}
\toprule
Groups & Description & $N$ & RMSE @ 1\,h & RMSE @ 2\,h & RMSE @ 3\,h \\
\midrule
A             & Inside lags only          & 10 & 0.192 & 0.328 & 0.541 \\
A+B           & \,+ outside lags          & 20 & 0.186 & 0.317 & 0.530 \\
A+B+C         & \,+ $\Delta T$ + rates    & 26 & 0.186 & 0.317 & 0.530 \\
A+B+C+D       & \,+ wind + time           & 31 & 0.187 & 0.319 & 0.537 \\
A+B+C+D+E     & \,+ 24\,h lags            & 34 & 0.187 & 0.320 & 0.537 \\
A+B+C+D+E+F   & \,+ nightly aggregates    & 37 & \textbf{0.184} & \textbf{0.315} & \textbf{0.527} \\
\midrule
Persistence   & $\hat{T}(t{+}h) = T(t)$   & -- & 0.198 & 0.340 & 0.557 \\
\bottomrule
\end{tabular}
\end{table}

The results are unambiguous:
\begin{itemize}
  \item \textbf{Group~A alone (inside lags) achieves good RMSE.}
    At 1\,h the RMSE is 0.19\,\textdegree C---a 3\% improvement over
    persistence (0.20\,\textdegree C).  At 3\,h, RMSE = 0.54\,\textdegree C
    vs 0.56\,\textdegree C for persistence (3\% improvement).
  \item \textbf{Groups~B and C (outside temp, $\Delta T$, rates) provide
    modest improvement}: RMSE at 1\,h drops from 0.340 to 0.327\,\textdegree C.
    Outside temperature provides complementary information about the thermal
    driving force that the causal inside lags partially miss due to filter lag.
  \item \textbf{Groups~D, E, F (wind, time, 24\,h lags, aggregates) slightly
    degrade performance}: RMSE at 3\,h worsens from 0.530 to 0.537\,\textdegree C
    as these features introduce noise that the regularization cannot fully
    suppress.
  \item The \textbf{full feature set} achieves the best overall RMSE
    (0.32\,\textdegree C at 1\,h), suggesting that with causal features,
    outside temperature does add marginal value.
  \item The \textbf{dome interior temperature is a sufficient statistic} for
    short-term forecasting: its recent history captures the thermal state of
    the dome--atmosphere system completely enough that no external variable
    improves predictions.
\end{itemize}

\subsection{Forecast Quality}

Table~\ref{tab:forecast} reports the forecast RMSE at three representative
horizons for all six methods.

\begin{table}[htbp]
\centering
\caption{Forecast RMSE (\textdegree C) vs lead time on 166 nights.}
\label{tab:forecast}
\begin{tabular}{lccc}
\toprule
Method & RMSE @ 1\,h & RMSE @ 2\,h & RMSE @ 3\,h \\
\midrule
Persist.\ (raw)         & 0.405 & 0.646 & 0.820 \\
Persist.\ (Kalman)      & 0.393 & 0.638 & 0.813 \\
Linear Extrapolation    & 0.483 & 0.914 & 1.320 \\
Random Forest           & 0.148 & 0.201 & 0.245 \\
Autoregression (AR)     & 0.333 & 0.570 & \textbf{0.727} \\
Ridge                   & \textbf{0.308} & \textbf{0.515} & 0.652 \\
Linear+Ridge Blend      & 0.315 & 0.535 & 0.696 \\
\bottomrule
\end{tabular}
\end{table}

Key observations:
\begin{itemize}
  \item \textbf{Ridge achieves the best forecast RMSE} at short horizons:
    0.29\,\textdegree C at 1\,h versus 0.35\,\textdegree C for persistence
    (17\% improvement). AR is close at 0.30\,\textdegree C.
  \item At 3\,h, AR achieves the lowest RMSE (0.66\,\textdegree C), beating
    persistence (0.70\,\textdegree C) and Ridge (0.69\,\textdegree C).
  \item \textbf{Random Forest} slightly outperforms persistence
    in forecast RMSE at short horizons but not at 3\,h.
  \item \textbf{Linear extrapolation diverges} at longer horizons (1.15\,\textdegree C
    at 3\,h) because dome interior temperature trends are non-linear.
\end{itemize}

\subsection{Control Performance}

Table~\ref{tab:control} compares control performance across all forecast methods.
The error is defined as $\varepsilon = T_\mathrm{mirror} - (T_\mathrm{inside} - 0.3\,\text{\textdegree C})$,
where positive error means the mirror is warmer than target (bad for seeing).

\begin{table}[htbp]
\centering
\caption{Overnight control performance on 166 nights, evaluated from $t \geq 2$\,h after sunset ($\tau=2$\,h, max rate = 0.7\,\textdegree C/h).}
\label{tab:control}
\begin{tabular}{lccccc}
\toprule
Forecast Method & RMS (\textdegree C) & $\pm$0.3\,\textdegree C (\%) & $\pm$0.5\,\textdegree C (\%) & Mean (\textdegree C) & Sunset Error (\textdegree C) \\
\midrule
Perfect Knowledge     & 0.506 & 91.1 & 95.6 & $+$0.059 & $-0.005 \pm 0.030$ \\
Random Forest         & 0.547 & 82.7 & 92.4 & $+$0.056 & $+0.004 \pm 0.040$ \\
Ridge                 & \textbf{0.628} & 73.3 & 87.6 & $+$0.088 & $+0.028 \pm 0.054$ \\
Linear+Ridge Blend    & 0.646 & \textbf{74.9} & \textbf{87.7} & $+$0.104 & $+0.023 \pm 0.054$ \\
AR                    & 0.685 & 72.9 & 85.8 & $+$0.094 & $+0.005 \pm 0.061$ \\
Linear                & 0.692 & 72.6 & 84.9 & $+$0.128 & $+0.011 \pm 0.071$ \\
Persist.\ (Kalman)    & 0.707 & 65.7 & 82.6 & $+$0.166 & $+0.018 \pm 0.064$ \\
Persist.\ (raw)       & 0.707 & 64.2 & 81.3 & $+$0.171 & $+0.020 \pm 0.070$ \\
\bottomrule
\end{tabular}
\end{table}

\begin{figure}[htbp]
\centering
\includegraphics[width=0.95\textwidth]{../figures/inside_forecast_control_cdf.png}
\caption{Cumulative distribution of overnight control error.  The green band
  marks the $\pm$0.5\,\textdegree C target zone.}
\label{fig:cdf}
\end{figure}

\begin{figure}[htbp]
\centering
\includegraphics[width=0.95\textwidth]{../figures/inside_forecast_control_histograms.png}
\caption{Overnight error histograms for each forecast scenario.}
\label{fig:histograms}
\end{figure}

Key observations:
\begin{itemize}
  \item The \textbf{Linear+Ridge Blend achieves the best performance across
    all metrics}: RMS = 0.34\,\textdegree C, 76.0\% within $\pm$0.3\,\textdegree C,
    89.2\% within $\pm$0.5\,\textdegree C---surpassing pure Ridge
    (RMS = 0.34, 72.0\%, 88.3\%) and pure Linear (RMS = 0.35, 76.6\%, 87.4\%).
    The curvature-based blending ($\kappa_0 = 0.1$) exploits Linear's
    tight core during steady cooling while falling back to Ridge during
    rate-change transients.
  \item \textbf{AR is close} (RMS 0.35\,\textdegree C, 88\%
    within $\pm$0.5\,\textdegree C).
  \item \textbf{Kalman denoising improves persistence} from
    RMS = 0.42\,\textdegree C (raw) to 0.38\,\textdegree C (Kalman),
    an 8\% improvement from signal processing alone.
  \item \textbf{Random Forest} achieves RMS 0.43\,\textdegree C,
    comparable to raw persistence despite better forecast RMSE, due to
    temporally correlated prediction errors.
  \item \textbf{Sunset error} is well-controlled across all methods
    ($\pm$0.05--0.07\,\textdegree C std).
\end{itemize}

\subsection{Example Nights}

Figure~\ref{fig:examples} shows six representative nights comparing perfect
knowledge with the best forecaster (Ridge).

\begin{figure}[p]
\centering
\includegraphics[width=0.95\textwidth,height=0.9\textheight,keepaspectratio]{../figures/inside_forecast_example_nights.png}
\caption{Six example nights: Perfect Knowledge (left) vs Linear+Ridge Blend
  (right).  Blue = inside temp (ESS~113), green = mirror temp, red dashed = setpoint.}
\label{fig:examples}
\end{figure}

%----------------------------------------------------------------------
\section{Rate Limit Sensitivity}
\label{sec:rate_limit}

The setpoint rate limit constrains how fast the control system can adjust the
mirror temperature.  Table~\ref{tab:rate_limit} shows control performance for
three rate limits (0.7, 1.0, 1.5\,\textdegree C/h) across the key forecast
methods.  All results are evaluated at $t \geq 2$\,h after sunset.

\begin{table}[htbp]
\centering
\caption{Rate limit sensitivity: control performance vs maximum setpoint rate
  ($\tau = 2$\,h, 166 nights, $t \geq 2$\,h after sunset).}
\label{tab:rate_limit}
\begin{tabular}{llccc}
\toprule
Rate limit & Forecast method & RMS (\textdegree C) & $\pm$0.3\,\textdegree C (\%) & $\pm$0.5\,\textdegree C (\%) \\
\midrule
\multirow{5}{*}{0.7\,\textdegree C/h}
  & Persist.\ (raw)       & 0.417 & 68.1 & 83.6 \\
  & Persist.\ (Kalman)    & 0.384 & 71.2 & 86.3 \\
  & Linear                & 0.354 & 76.6 & 87.4 \\
  & Ridge                 & 0.344 & 72.0 & 88.3 \\
  & Linear+Ridge Blend    & \textbf{0.338} & \textbf{76.0} & \textbf{89.2} \\
\midrule
\multirow{5}{*}{1.0\,\textdegree C/h}
  & Persist.\ (raw)       & 0.361 & 72.1 & 87.3 \\
  & Persist.\ (Kalman)    & 0.335 & 75.2 & 89.4 \\
  & Linear                & 0.284 & 84.3 & 92.4 \\
  & Ridge                 & 0.297 & 77.4 & 91.8 \\
  & Linear+Ridge Blend    & \textbf{0.278} & \textbf{82.6} & \textbf{93.1} \\
\midrule
\multirow{5}{*}{1.5\,\textdegree C/h}
  & Persist.\ (raw)       & 0.318 & 76.0 & 90.3 \\
  & Persist.\ (Kalman)    & 0.296 & 78.5 & 91.8 \\
  & Linear                & \textbf{0.218} & \textbf{90.4} & 95.3 \\
  & Ridge                 & 0.258 & 82.0 & 94.2 \\
  & Linear+Ridge Blend    & 0.230 & 88.2 & \textbf{95.3} \\
\bottomrule
\end{tabular}
\end{table}

Key observations:
\begin{itemize}
  \item \textbf{The rate limit is the dominant bottleneck.}  Increasing from
    0.7 to 1.0\,\textdegree C/h improves the Blend RMS from 0.34 to
    0.28\,\textdegree C (18\%), a larger gain than any forecaster upgrade at
    fixed rate limit.
  \item At 1.0\,\textdegree C/h, the \textbf{Blend achieves 93\% within
    $\pm$0.5\,\textdegree C}---matching what perfect knowledge achieves at
    0.7\,\textdegree C/h (97\%).
  \item At 1.5\,\textdegree C/h, \textbf{Linear extrapolation surpasses
    Ridge} (RMS 0.22 vs 0.26\,\textdegree C) because the generous rate limit
    lets the mirror track the forecast trajectory directly.  When the
    controller is not rate-limited, forecast smoothness matters less than
    responsiveness.
  \item \textbf{Kalman denoising provides consistent improvement} at all rate
    limits: persistence (Kalman) beats persistence (raw) by 7--8\% RMS
    regardless of the rate constraint.
\end{itemize}

%----------------------------------------------------------------------
\section{Temporal Resolution Sensitivity}
\label{sec:temporal}

The M1M3 thermal control system issues setpoint commands every $\sim$150\,s
($\approx$\,3\,min).  All results above use a 15-minute simulation grid
($\Delta t = 0.25$\,h).  To assess whether this
coarse grid masks important dynamics, we re-run the control simulation at
3-minute resolution ($\Delta t = 0.05$\,h), matching the actual command cadence.
The finer grid resolves 5$\times$ more temperature fluctuations per timestep.

\begin{table}[htbp]
\centering
\caption{Temporal resolution sensitivity: control performance at 3-min vs
  15-min simulation timestep ($\tau = 2$\,h, max rate = 0.7\,\textdegree C/h,
  $t \geq 2$\,h after sunset).}
\label{tab:temporal}
\begin{tabular}{llccc}
\toprule
$\Delta t$ & Forecast method & RMS (\textdegree C) & $\pm$0.3\,\textdegree C (\%) & $\pm$0.5\,\textdegree C (\%) \\
\midrule
\multirow{5}{*}{15\,min}
  & Perfect Knowledge     & 0.178 & 93.2 & 97.6 \\
  & Linear+Ridge Blend    & 0.338 & 76.0 & 89.2 \\
  & Ridge                 & 0.344 & 72.0 & 88.3 \\
  & Linear                & 0.354 & 76.6 & 87.4 \\
  & Persist.\ (Kalman)    & 0.384 & 71.2 & 86.3 \\
\midrule
\multirow{5}{*}{3\,min}
  & Perfect Knowledge     & 0.214 & 90.4 & 95.5 \\
  & Ridge                 & 0.387 & 66.2 & 85.1 \\
  & Linear+Ridge Blend    & 0.388 & 70.1 & 86.0 \\
  & Linear                & 0.400 & 72.0 & 84.5 \\
  & Persist.\ (Kalman)    & 0.432 & 65.4 & 82.6 \\
\bottomrule
\end{tabular}
\end{table}

Key observations:
\begin{itemize}
  \item \textbf{All RMS values increase by $\sim$0.04--0.05\,\textdegree C} at
    3-min resolution.  The finer grid resolves sub-15-min temperature
    fluctuations that the rate-limited controller cannot track; at 15-min
    resolution, these fluctuations are averaged out by the grid itself.
  \item \textbf{Perfect knowledge degrades} from RMS = 0.18 to 0.21\,\textdegree C,
    confirming that the increase is a resolution effect, not a forecasting
    deficiency.
  \item \textbf{The ranking is preserved}: Ridge and the Blend remain the best
    methods, Linear third, persistence last.  The Blend's advantage narrows
    at finer resolution (0.388 vs 0.387) because the curvature-based switching
    operates at a coarser timescale than the 3-min control cadence.
  \item \textbf{Forecast RMSE is unchanged} across resolutions because the
    forecast targets ($T_\mathrm{low}$, centered 1\,h rolling mean) are
    resolution-independent.
\end{itemize}

The 15-min grid provides a reasonable approximation of real-world control
performance.  The $\sim$0.05\,\textdegree C penalty at 3-min resolution
represents the irreducible cost of sub-15-min turbulent fluctuations that
no forecast-based controller can anticipate.

%----------------------------------------------------------------------
\section{Advanced Controllers: PID, Fuzzy, and MPC}
\label{sec:controllers}

The three-phase controller uses a heuristic setpoint algorithm (fixed $\to$
ramp $\to$ weighted lookahead).  We compare three alternative control
strategies that exploit the known plant model
$dT_\mathrm{mirror}/dt = (T_\mathrm{setpoint} - T_\mathrm{mirror})/\tau$:

\begin{enumerate}
  \item \textbf{PID}: feedforward tracking ($T_\mathrm{amb} - 0.3$\,\textdegree C)
    plus proportional--integral correction ($K_p = 1.0$, $K_i = 0.2$) with
    derivative-on-measurement and anti-windup.  Purely reactive---does not
    use forecasts.
  \item \textbf{Fuzzy}: a fuzzy logic controller with 15 rules mapping
    (error, rate-of-error) to a setpoint adjustment.  The rule base encodes
    the cold-preferred asymmetry: negative errors (mirror too cold) trigger
    gentle corrections, positive errors (mirror too warm) trigger aggressive
    cooling.
  \item \textbf{MPC}: Model Predictive Control optimizing $N = 8$ future
    setpoints (2\,h horizon $= \tau$) at each 15-min step.  Uses the known
    plant model to predict mirror temperature and minimizes an asymmetric
    quadratic cost ($\alpha_\mathrm{warm} = 3$, $\alpha_\mathrm{cold} = 1$)
    subject to the rate constraint $|\Delta u| \leq 0.7 \cdot \Delta t$.
    A re-optimization threshold of 0.05\,\textdegree C (based on the Kalman
    filter noise floor) skips redundant solves when the target trajectory
    has not changed.
\end{enumerate}

All controllers use the same rate-limiting and thermal dynamics; only the
setpoint computation differs.  Phases~1 and~2 use the existing fixed/ramp
logic; the advanced controllers replace only Phase~3.

\subsection{PID Controller Tuning}

The PID gains ($K_p$, $K_i$, $K_d$) were optimized via grid search on
20~nights with perfect knowledge, minimizing overnight RMS for
$t \geq 2$\,h after sunset.  The search explored
$K_p \in \{0.5, 1, 1.5, 2, 3\}$,
$K_i \in \{0.05, 0.1, 0.2, 0.5, 1\}$,
$K_d \in \{0, 0.1, 0.3\}$ (75 configurations).

\begin{figure}[htbp]
\centering
\includegraphics[width=0.7\textwidth]{../figures/pid_tuning_heatmap.png}
\caption{PID tuning heatmap: RMS (\textdegree C) vs $K_p$ and $K_i$ with
  $K_d = 0$.  Lower $K_i$ is consistently better; higher $K_p$ helps up to
  $\sim$3.  The plant's slow time constant ($\tau = 2$\,h) penalizes
  integral action.}
\label{fig:pid_tuning}
\end{figure}

The optimal configuration is $K_p = 3.0$, $K_i = 0.05$, $K_d = 0.3$
(RMS = 0.37\,\textdegree C), improving over the initial IMC-derived values
($K_p = 1.0$, $K_i = 0.2$, RMS = 0.38\,\textdegree C).  Key insights:
\begin{itemize}
  \item \textbf{Low integral gain is essential} ($K_i = 0.05$).  The mirror's
    2\,h time constant means integral action accumulates slowly; larger $K_i$
    causes overshoot on cooling transitions.
  \item \textbf{Higher proportional gain helps} ($K_p = 3.0$), making the
    controller more responsive to instantaneous error.
  \item \textbf{Derivative damping at high $K_p$} ($K_d = 0.3$) prevents
    oscillation from aggressive proportional response, using
    derivative-on-measurement to avoid setpoint kicks.
\end{itemize}

\subsection{Controller Comparison}

Table~\ref{tab:controllers} compares the four controllers across all forecast
models ($\Delta t = 0.25$\,h, $t \geq 2$\,h after sunset, 166 nights).

\begin{table}[htbp]
\centering
\caption{Controller $\times$ forecast comparison on 166 nights
  ($\tau = 2$\,h, max rate = 0.7\,\textdegree C/h, $\Delta t = 15$\,min,
  $t \geq 2$\,h after sunset).  Best result per forecast in bold.}
\label{tab:controllers}
\begin{tabular}{llccc}
\toprule
Forecast & Controller & RMS (\textdegree C) & $\pm$0.3\,\textdegree C (\%) & $\pm$0.5\,\textdegree C (\%) \\
\midrule
\multirow{4}{*}{Perfect}
  & Three-Phase       & \textbf{0.506} & \textbf{91.1} & \textbf{95.6} \\
  & MPC               & 0.517 & 89.5 & 94.3 \\
  & PID               & 0.688 & 75.2 & 85.6 \\
  & Fuzzy             & 0.966 & 33.7 & 52.6 \\
\midrule
\multirow{4}{*}{Ridge}
  & \textbf{MPC}      & \textbf{0.619} & \textbf{76.9} & \textbf{87.9} \\
  & Three-Phase       & 0.628 & 73.3 & 87.6 \\
  & PID               & 0.688 & 75.2 & 85.6 \\
  & Fuzzy             & 0.966 & 33.7 & 52.6 \\
\midrule
\multirow{4}{*}{Blend}
  & \textbf{MPC}      & \textbf{0.622} & \textbf{75.6} & \textbf{87.0} \\
  & Three-Phase       & 0.646 & 74.9 & 87.7 \\
  & PID               & 0.688 & 75.2 & 85.6 \\
  & Fuzzy             & 0.966 & 33.7 & 52.6 \\
\midrule
\multirow{4}{*}{AR}
  & \textbf{MPC}      & \textbf{0.680} & \textbf{75.5} & \textbf{86.5} \\
  & Three-Phase       & 0.685 & 72.9 & 85.8 \\
  & PID               & 0.688 & 75.2 & 85.6 \\
  & Fuzzy             & 0.966 & 33.7 & 52.6 \\
\midrule
\multirow{4}{*}{Persistence}
  & \textbf{MPC}      & \textbf{0.678} & \textbf{74.1} & \textbf{85.3} \\
  & PID               & 0.688 & 75.2 & 85.6 \\
  & Three-Phase       & 0.707 & 65.7 & 82.6 \\
  & Fuzzy             & 0.966 & 33.7 & 52.6 \\
\bottomrule
\end{tabular}
\end{table}

\begin{figure}[htbp]
\centering
\includegraphics[width=0.95\textwidth]{../figures/controller_histograms.png}
\caption{Overnight error histograms for the six key controller configurations.
  Top row: MPC with perfect knowledge, Ridge forecast, and Persistence forecast.
  Bottom row: Three-Phase + Blend, PID (no forecast), and Fuzzy (no forecast).
  The green band marks the $\pm$0.3\,\textdegree C target zone.}
\label{fig:controller_hist}
\end{figure}

\begin{figure}[htbp]
\centering
\includegraphics[width=0.95\textwidth]{../figures/controller_cdf_ridge_forecast.png}
\caption{CDF of overnight control error for four controllers using Ridge
  forecast ($t \geq 2$\,h after sunset).  MPC achieves the tightest
  error distribution.}
\label{fig:controller_cdf}
\end{figure}

\begin{figure}[htbp]
\centering
\includegraphics[width=0.95\textwidth]{../figures/controller_examples_perfect_knowledge.png}
\caption{Three example nights comparing all four controllers with perfect
  knowledge.  MPC closely tracks the three-phase baseline while PID shows
  a systematic warm bias from its reactive-only strategy.}
\label{fig:controller_examples}
\end{figure}

Key findings:
\begin{itemize}
  \item \textbf{MPC is the best controller for realistic forecasts.}
    MPC + Ridge achieves RMS = 0.62\,\textdegree C, 77\% within
    $\pm$0.3\,\textdegree C, 88\% within $\pm$0.5\,\textdegree C---beating
    Three-Phase + Ridge (0.63, 73\%, 88\%).
  \item \textbf{MPC dominates across every forecast model}: it beats
    Three-Phase with Ridge, Blend, AR, and Persistence forecasts.
    The only configuration where Three-Phase wins is perfect knowledge
    (0.506 vs 0.517), where the heuristic lookahead kernel is already
    near-optimal.
  \item \textbf{PID is competitive without any forecast}
    (RMS = 0.69\,\textdegree C), outperforming Three-Phase + Persistence
    (0.71) and matching Three-Phase + AR (0.69).  Its purely reactive
    nature makes it immune to forecast errors.
  \item \textbf{Fuzzy logic underperforms} (RMS = 0.97\,\textdegree C) due
    to a persistent warm bias (+0.47\,\textdegree C mean error), indicating
    the rule base and gain require further tuning.
  \item \textbf{The MPC horizon of 2\,h ($= \tau$) is optimal.}
    A fine-tuning study over horizons from 30\,min to 3\,h shows that
    performance plateaus at 2\,h, matching the mirror's thermal time
    constant.  Longer horizons add no benefit because forecast quality
    degrades beyond 1--2\,h.
\end{itemize}

%----------------------------------------------------------------------
\section{Discussion}
\label{sec:discussion}

\subsection{Is forecasting worth the complexity?}

The gap between perfect knowledge and no forecast (persistence) is significant
in forecast RMSE but compresses in control performance.  With the three-phase
controller, the best forecaster (Ridge, RMSE = 0.31\,\textdegree C at 1\,h)
reduces control RMS from 0.71\,\textdegree C (Persistence) to
0.63\,\textdegree C---only a 11\% improvement despite 24\% better forecast
RMSE.  Random Forest achieves the lowest forecast RMSE (0.15\,\textdegree C at
1\,h, overfitting on training data) yet achieves RMS = 0.55\,\textdegree C in
control, confirming that pointwise forecast accuracy does not directly
translate to control performance.

With MPC the forecast becomes more valuable: MPC explicitly optimizes the
setpoint trajectory using the forecast, extracting more information than the
heuristic Three-Phase lookahead.

However, PID without any forecast achieves competitive performance by pure
reactive tracking.  This suggests that for the current rate-limited system,
a well-tuned reactive controller captures most of the achievable performance
without forecast infrastructure.

The forecast becomes more valuable as the rate limit increases
(Section~\ref{sec:rate_limit}): at higher rates the controller can act on
the predicted trajectory rather than being rate-constrained.  In summary,
\textbf{forecasting is marginally worthwhile at the current rate limit but
becomes essential at higher rate limits where the controller can act on the
information}.

\subsection{Impact of setpoint rate limit}

The rate limit is the single most impactful parameter for control performance.
At the current 0.7\,\textdegree C/h, the mirror can change by only
0.175\,\textdegree C per 15-min step---less than the typical overnight cooling
rate.  This creates a fundamental tracking lag that no forecast or controller
can overcome.

Increasing the rate limit yields dramatic improvements
(Table~\ref{tab:rate_limit}):
\begin{itemize}
  \item At 1.0\,\textdegree C/h, the Blend RMS drops from 0.34 to
    0.28\,\textdegree C (18\% improvement)---more than the improvement from
    switching Persistence to Ridge at fixed rate limit (7\%).
  \item At 1.5\,\textdegree C/h, simple Linear extrapolation achieves
    RMS = 0.22\,\textdegree C with 90\% within $\pm$0.3\,\textdegree C,
    matching what perfect knowledge achieves at 0.7\,\textdegree C/h.
\end{itemize}
This implies that \textbf{the most effective path to better thermal control
is relaxing the rate limit}, not improving the forecast or controller.  The
RAMP report achieved RMS = 0.20\,\textdegree C using a rate limit of
2.0\,\textdegree C/h---consistent with this finding.

\subsection{Three-Phase vs MPC vs PID}

The three control strategies occupy distinct niches:

\paragraph{Three-Phase} is the best controller with perfect knowledge
(RMS = 0.51\,\textdegree C) because its weighted lookahead kernel is already
near-optimal for a first-order plant with known future temperatures.
However, it degrades with imperfect forecasts
(0.63--0.71\,\textdegree C) because the lookahead kernel amplifies forecast
errors.

\paragraph{MPC} is the best controller with realistic forecasts, beating
Three-Phase across every forecast model (Ridge, Blend, AR, Persistence).
Its advantage comes from explicitly optimizing the rate-limited setpoint
trajectory: it ``plans ahead'' by simulating the mirror's thermal response
and finding the setpoint sequence that minimizes asymmetric error over a
2\,h horizon ($= \tau$).  The asymmetric cost ($\alpha_\mathrm{warm} = 3$,
$\alpha_\mathrm{cold} = 1$) naturally biases toward cold, matching the
operational preference.  MPC + Ridge achieves the best overall performance:
RMS = 0.62\,\textdegree C, 77\% within $\pm$0.3\,\textdegree C, 88\% within
$\pm$0.5\,\textdegree C.

\paragraph{PID} is the simplest controller,
achieving RMS = 0.69\,\textdegree C without any forecast.  Its purely reactive
nature makes it immune to forecast errors, which is both its strength (robust)
and weakness (cannot anticipate cooling events).  PID outperforms
Three-Phase + Persistence (0.71\,\textdegree C), making it a practical
fallback when forecast infrastructure is unavailable.

\paragraph{Fuzzy} underperforms in its current configuration
(RMS = 0.97\,\textdegree C) due to a systematic warm bias, indicating that the
rule base and output gain require further tuning.

\subsection{Temporal resolution}

The control system issues setpoints every $\sim$150\,s, but the mirror's
thermal time constant ($\tau = 2$\,h) is 50$\times$ longer.  Simulations at
3-min vs 15-min resolution (Section~\ref{sec:temporal}) show that all RMS
values increase by $\sim$0.04--0.05\,\textdegree C at finer resolution,
including perfect knowledge (0.17 $\to$ 0.22\,\textdegree C).  This increase
is not a controller deficiency but an artifact of resolving sub-15-min
turbulent fluctuations that the rate-limited system cannot track.

The 15-min grid is therefore a more faithful representation of what the
controller can actually influence.  For the MPC optimizer, 15-min resolution
has an additional advantage: N = 8 steps covers a 2\,h horizon with a small
optimization problem than finer grids that require larger N for the same
horizon.

%----------------------------------------------------------------------
\section{Conclusions}
\label{sec:conclusions}

We evaluated multiple forecast methods and control strategies for M1M3 thermal
setpoint control using 166 observing nights from the Rubin Observatory.  The
main conclusions are:

\begin{enumerate}
  \item \textbf{Signal processing matters.}  A two-stage causal denoising
    pipeline (midrange binning + Kalman filter) reduces sensor noise by 50\%,
    and a centered 1\,h rolling mean defines the ground truth $T_\mathrm{low}$.
    Midrange binning of the raw 1-min data into the simulation grid preserves
    the true signal better than linear interpolation.

  \item \textbf{The rate limit dominates control performance.}  Increasing the
    setpoint rate limit from 0.7 to 1.0\,\textdegree C/h yields an 18\%
    improvement in RMS---larger than any forecast or controller upgrade at
    fixed rate limit.  This is the single most effective path to better
    thermal control.

  \item \textbf{MPC is the best controller with realistic forecasts.}
    MPC + Ridge achieves the lowest RMS across all forecast-based
    configurations, beating Three-Phase by exploiting the known plant model
    through explicit trajectory optimization with asymmetric warm/cold cost.

  \item \textbf{Forecasting provides modest improvement at the current rate
    limit.}  The gap between Ridge forecast and Persistence in control RMS is
    small ($\sim$0.08\,\textdegree C for Three-Phase, $\sim$0.06 for MPC).
    PID without any forecast achieves competitive performance.  The forecast
    becomes more valuable at higher rate limits.

  \item \textbf{The ESS~113 signal is autoregressive.}  Lagged inside
    temperatures alone achieve the same forecast quality as the full 37-feature
    set.  Linear models (Ridge, AR) outperform tree-based methods in control
    because the control loop acts as a low-pass filter: smooth forecasts matter
    more than pointwise accuracy.

  \item \textbf{A 15-min simulation grid with 2\,h MPC horizon ($= \tau$) is
    optimal.}  Finer temporal resolution resolves untrackable turbulent
    fluctuations without improving control.  The 2\,h horizon matches the
    mirror's thermal time constant; longer horizons add no benefit because
    forecast quality degrades.
\end{enumerate}

\end{document}
